\section{Estimation}
\label{sec:estimation}

\subsection{State-space form}

Let $\bm{x}_t$ be the vector of predetermined variables of the linearized model (capital and the exogenous processes), with law of motion $\bm{x}_{t+1} = \bm{T}\bm{x}_t + \bm{R}\bm{\varepsilon}_{t+1}$, and let all model variables be $\bm{M}\bm{x}_t$. Each observable in \eqref{eq:obs} is a difference of model variables plus the trend shock for the variables with a trend, for example $\Delta\ln Y_t = \hat y_t - \hat y_{t-1} + \mu_t$. The state is therefore augmented with its lag.

Survey and statistical noise is modeled in \emph{levels}. For investment, employment and the logit informality rate, the observed log level equals the model level plus an independent error $e_{j,t}\sim N(0,\sigma^2_{e_j})$, so that the growth rate contains $e_{j,t}-e_{j,t-1}$. These errors are added as states. This specification produces the negative first-order autocorrelation of the corresponding growth rates and lets the Kalman smoother separate survey noise from movements in the latent informality rate. A small error with standard deviation $10^{-3}$ is added to every growth rate to keep the innovation covariance non-singular. The likelihood is evaluated with the Kalman filter, initialized at the unconditional distribution; quarters with missing observations skip the update step \citep{DurbinKoopman2012}. The implementation was checked against the exact Gaussian likelihood of the stacked observations. Appendix~\ref{app:model} writes every measurement equation explicitly and explains the filter and the sampler step by step.

\subsection{Calibration}

Table~\ref{tab:calibration} lists the parameters fixed before estimation. Those taken from the data are sample means of 1993--2026. The capital share and depreciation come from Colombian sources and from the data, as described below; the equal-technology estimates use $\alpha=0.33$, for which the depreciation implied by the investment ratio is also 4\% a year. The informal share $\theta$ is \emph{not} calibrated: it is estimated, and its steady state is recomputed at every draw. Section~\ref{sec:identification} shows that the object of interest, the path of the informal share, does not depend on its value.

\begin{table}[H]
\centering
\caption{Calibrated parameters.}
\label{tab:calibration}
\footnotesize
\begin{tabularx}{\textwidth}{L{1.6cm} L{6.2cm} r Y}
\toprule
\textbf{Parameter} & \textbf{Meaning} & \textbf{Value} & \textbf{Source} \\
\midrule
$\alpha$ & Capital share of formal output & 0.354 & Estimated, Section~\ref{sec:alpha} \\
$\beta$ & Quarterly discount factor & 0.985 & Conventional, emerging economy \\
$\varphi$ & Inverse Frisch elasticity & 1.0 & Conventional \\
$\bar\tau$ & Steady-state formality wedge & 0.25 & Affects only the scale of $\chi_f$ \\
$\gamma$ & Mean quarterly trend growth & 0.0096 & GDP, 2005--2018 \\
$\bar s$ & Mean informality rate & 0.645 & ICML, 1993--2026 \\
$C/Y$, $I/Y$, $G/Y$ & Expenditure ratios & 0.697, 0.186, 0.139 & National accounts, 1993--2026 \\
$\delta$ & Quarterly depreciation & 0.0102 & 4\% a year; GRECO 3.1\%, DNP 3.8\% \\
$\bar x$ & Chain-volume residual & -- & Implied by expenditure ratios \\
$\ln(1+d_t^{\mathrm{nw}})$ & Non-wage labor costs of formal employees & 0.42 / 0.32 & Statutory rates, 2008--2012 / 2014--; Section~\ref{sec:data} \\
\bottomrule
\end{tabularx}
\end{table}

\subsection{Capital share and depreciation}
\label{sec:alpha}

Depreciation is set at the standard 4\% a year. Colombian estimates are close: the capital stocks of the Banco de la Rep\'ublica growth project (GRECO), built with the method of \citet{Barrios1993}, imply 2\% for buildings, 5\% for machinery and equipment and 3.1\% for total fixed capital between 1950 and 1997 \citep{GRECO1999}, and the National Planning Department uses 3.8\% \citep{RodriguezPerillaReyes2004}.

The capital share of formal output is estimated. Capital is built by the perpetual-inventory method from gross fixed capital formation in constant prices, and
\begin{equation}
\ln\frac{Y_{f,t}}{N_{f,t}} = c + \alpha\ln\frac{K_t}{N_{f,t}} + g\,t + e_t
\label{eq:dols}
\end{equation}
is estimated by dynamic OLS \citep{StockWatson1993} with one lead and lag and Newey--West errors on annual data for 1993--2025. Formal output $Y_{f,t}=(1-\theta_t)Y_t$ and formal employment $N_{f,t}=(1-s_t)N_t$ depend on $\alpha$ through \eqref{eq:share}, so $\alpha$ is a fixed point. With 4\% depreciation the estimate is 0.354 (standard error 0.25); it is 0.33 with 6\% and 0.30 with 8\%. The regression without trend gives 0.39--0.43, which absorbs productivity growth into capital, and the aggregate regression gives an implausible 0.65. The time series is short and the estimate is imprecise, but it is in the range of earlier Colombian estimates: GRECO finds 0.42 for 1925--1981 by cointegration and 0.36--0.41 for 1926--1994 and 1950--1994 by restricted least squares \citep{GRECO1999}, and the National Planning Department 0.39--0.51 depending on the specification \citep{RodriguezPerillaReyes2004}. Sensitivities with $\alpha=0.33$ (and 6\% depreciation) and $\alpha=0.50$ are reported in Section~\ref{sec:results}.

\subsection{Priors and sampler}

The estimated parameters are the informal share $\theta$, the labor elasticity of informal production $\eta$ (labor-only model), the seven persistences and seven innovation standard deviations (including the relative-productivity shock), the four measurement-error standard deviations (three level errors and the error of the earnings ratio) and the level shift $b$ of the earnings ratio from 2022Q1. Persistences, $\theta$ and $\eta$ have beta priors ($\eta$ with mean 0.7 and standard deviation 0.1), standard deviations inverse-gamma priors and $b$ a normal prior with mean zero and standard deviation 0.2 (Table~\ref{tab:posterior}). The prior of $\theta$, with mean 0.35 and standard deviation 0.10, is deliberately wide: its 90\% interval covers informal shares from about 19\% to 52\% of output.

The posterior is sampled in an unconstrained parameterization (logit for parameters in the unit interval, logarithm for standard deviations, identity for $b$), including the Jacobian. The posterior mode is found from three starting points by Nelder--Mead followed by BFGS; the inverse Hessian at the mode, computed by finite differences, is the proposal covariance of a random-walk Metropolis--Hastings sampler \citep{AnSchorfheide2007}. Eight chains start from over-dispersed points around the mode; during 15{,}000 burn-in iterations the proposal scale is adapted towards an acceptance rate of 0.28 \citep{RobertsGelmanGilks1997}, and 40{,}000 further iterations per chain are kept with a fixed proposal (four chains of 30{,}000 after 10{,}000 in the equal-technology and aggregates-only estimates). Convergence is assessed with the rank-normalized split $\hat R$ and the bulk and tail effective sample sizes of \citet{Vehtari2021}.
